\PassOptionsToPackage{dvipdfmx}{graphicx}
\documentclass[twocolumn, fontset = fandol]{ctexart}

\ExplSyntaxOn \makeatletter

\keys_define:nn { hduthesis / exam }
  {
    watermark .tl_set:N = \l__hduthesis_exam_watermark_tl,
  }
\NewDocumentCommand \hdutexset {m} { \keys_set:nn { hduthesis / exam } {#1} }
\DeclareHookRule { shipout / background }
  { hduthesis / exam } { before } { pgfrcs }
\AddToHook { shipout / background } [ hduthesis / exam ]
  {
    \tl_if_empty:NF \l__hduthesis_exam_watermark_tl
      {
        \tikz [ remember~picture, overlay ]
          \node [ opacity = .2 ] at (current~page)
            {
              \includegraphics [ width = \linewidth ]
                { \l__hduthesis_exam_watermark_tl }
            };
      }
  }

\RequirePackage
  {
    geometry, mathtools, array, bm, esint, keytheorems, libertine, inconsolata,
    siunitx, fixdif, derivative, physics2, tasks, anyfontsize
  }
\RequirePackage[inline]{enumitem}
\RequirePackage{fontawesome5, fancyqr, circuitikz}
\hook_gput_code:nnn { cmd/faIcon/before } { . }
  {
    \group_begin:
    \fontsize { \fp_eval:n { .9 * \f@size } } { 0 } \selectfont
  }
\hook_gput_code:nnn { cmd/faIcon/after } { . } { \group_end: }
\geometry{margin = 1in}

\RequirePackage[svgnames]{xcolor}
\newkeytheorem{note}
  [ style = remark, numbered = false,
    name = \color{DarkRed}\sffamily \faIcon{info-circle}~Tip. ]
\newkeytheorem{problem}
  [ style = definition,
    name = \color{DarkGreen}\sffamily \faIcon{pen-square}~Problem ]
\newkeytheorem{solution}
  [ style = definition, numbered = false,
    name = \color{MidnightBlue}\sffamily \faIcon{check-square}~Solution ]
\usephysicsmodule{ab, op.legacy}
\usephysicsmodule{bm-um.legacy}
\settasks { label = \sffamily\Alph*., item-indent = 2.5em, label-offset = .5pt }
\newcommand \true {\addtocounter{task}{1}\faIcon{check-circle}}
\def \@result#1 {\underline{~#1~}}
\def \s@result {\textsf{【\qquad 】}}
\DeclareRobustCommand \result {\@ifstar\s@result\@result}
\keys_define:nn { textfig }
  {
    ratio      .tl_set:N    = \l__textfig_ratio_tl,
      ratio    .initial:n   = 7 \c_colon_str 2,
    position   .tl_set:N    = \l__textfig_position_tl,
      position .initial:n   = c,
    unknown    .code:n      = \tl_if_novalue:nF {#1}
      { \tl_set_eq:NN \l__textfig_position_tl \l_keys_key_tl }
  }
\NewDocumentEnvironment{textfig}{ O{} }
  {
    \group_begin:
    \keys_set:nn { textfig } {#1}
    \seq_set_split:Nee \l__whats_textfig_ratio_seq { \c_colon_str }
      { \l__textfig_ratio_tl }
    \noindent
    \minipage [ \l__textfig_position_tl ]
      {
        \fp_eval:n
          {
            .96*\seq_item:Nn \l__whats_textfig_ratio_seq {1}/
              (
                \seq_item:Nn \l__whats_textfig_ratio_seq {1} + 
                \seq_item:Nn \l__whats_textfig_ratio_seq {2}
              )
          }\dimexpr \linewidth
      }
  }
  {
    \endminipage
    \vspace{4.5pt plus 1.5pt minus 1.5pt} \par
    \group_end:
  }
\NewDocumentCommand \textfigsplit {}
  {
    \endminipage
    \hspace{.02\linewidth}
    \minipage [ \l__textfig_position_tl ]
      {
        \fp_eval:n
          {
            .96*\seq_item:Nn \l__whats_textfig_ratio_seq {2}/
              (
                \seq_item:Nn \l__whats_textfig_ratio_seq {1} + 
                \seq_item:Nn \l__whats_textfig_ratio_seq {2}
              )
          }\dimexpr \linewidth
      }
  }

% \hook_gput_code:nnn { cmd/maketitle/before } { . } { \vskip -\baselineskip }
% \hook_gput_code:nnn { cmd/maketitle/after } { . }
%   { \__weyoung_maketitle_aftercode_aux: }
\cs_new_protected:Nn \__weyoung_maketitle_aftercode_aux:
  {
    \vspace*{-2\baselineskip}
    \begin{center}
      \begin{minipage}{.66\textwidth}
        \centering
        \begin{tasks}
          [ label-width = 5\ccwd, item-indent = 2cm,
            after-item-skip = 0pt ](2)
          \task[\textbf{课程信息}：] 理学院/A0715022
          \task[\textbf{授课教师}：] 陈许敏
          \task[\textbf{解析提供}：] \textsf{Skyrmion}
          \task[{\makebox[5em][s]{\textbf{\textsf\TeX~排版}：}}] \textsf{Skyrmion}
        \end{tasks}
      \end{minipage}
      \begin{minipage}{.15\textwidth}
        \centering
        \begin{tikzpicture}
          \node [ inner~sep = 0 ] (wechat) {
            \fancyqr
              [ size = 4em, image~padding = 1,
                image = \raisebox{4em/33}{\textcolor{DarkGreen}{\faIcon{weixin}}},
                r~color = DarkGreen, version = 4
              ]{http://weixin.qq.com/r/sRJiemTEGjllrSAZ90ej}
            };
          \node [ inner~sep = .5ex, align = center, right, font = \scriptsize ]
            at (wechat.east) {未\\央\\学\\社};
        \end{tikzpicture}
      \end{minipage}
      \begin{minipage}{.15\textwidth}
        \centering
        \begin{tikzpicture}
          \node [ inner~sep = 0 ] (qq) {
            \fancyqr
              [ size = 4em, image~x~padding = 2, image~y~padding = 1,
                image = \raisebox{4em/33}{\textcolor{MidnightBlue}{\faIcon{qq}}},
                l~color = MidnightBlue, version = 4
              ]{https://qm.qq.com/q/PH7CsQWvm2}
            };
          \node [ inner~sep = .5ex, align = center, left, font = \scriptsize ]
            at (qq.west) {物\\理\\营};
        \end{tikzpicture}
      \end{minipage}
    \end{center}
  }

\newcommand \iu {{\mathup i\mkern1mu}}
\newcommand \eu {\mathup e}
\AtBeginDocument{%
  \let \Re \relax \let \Im \relax
  \DeclareMathOperator{\Re}{Re}
  \DeclareMathOperator{\Im}{Im}
}

\makeatother \ExplSyntaxOff
\usetikzlibrary{ fadings, patterns, arrows, decorations.markings }
\tikzset
  {
    > = stealth',
    every picture/.append style =
      { line join = round, line cap = round, thick }
  }

% \setCJKmainfont{LXGW WenKai}[AutoFakeSlant]
% \setCJKsansfont{LXGW Marker Gothic}[AutoFakeBold, AutoFakeSlant, Scale = .95]
% \setCJKmonofont{LXGW WenKai Mono}[AutoFakeBold, AutoFakeSlant]
\ctexset{section/format+ = \raggedright}
\hdutexset{watermark = ../WY_Logo.pdf}
% \usepackage[warnings-off={mathtools-colon,mathtools-overbracket}]{unicode-math}
% \setmathfont[AutoFakeBold]{STIX Two Math}

\makeatletter \def \@result #1{\underline{ \quad }} \makeatother

\ctikzset
  {
    amplifiers/fill  = cyan!20,
    resistors/fill   = violet!20,
    sources/fill     = yellow!20,
    diodes/fill      = red!20,
    resistors/scale  = .64,
    capacitors/scale = .6,
    diodes/scale     = .6,
    sources/scale    = .8,
    inductors/scale  = .6,
  }
\tikzset {ev fonery node/.style = {font = \small, align = left}}

\title{\sffamily\bfseries 24Fall \textsc{Hdu} 「模拟电子技术」补考试题}
\author{} \date{}

\begin{document}

\maketitle

\begin{problem}[12 pt]% 1
  如图所示为加减运算电路，求输出电压 $v_o$ 的表达式.
  \begin{center}
    \begin{circuitikz}
      \draw (0,0) node [ op amp, anchor = out ] (A) {$\mathrm A$}
       to [ short, -o ] ++ (1,0) node [ right ] {$v_o$};
      \draw (A.-) --++ (-.5,0)
       to [ short, -o, european resistor, l_ = $R_1$, name = R1 ] ++
       (-2,0) node [ font = \footnotesize\ttfamily ] at (R1) {10k}
       node [ left ] {$v_{i_1}$};
      \draw (A.+) --++ (-.5,0)
       to [ short, -o, european resistor, l_ = $R_3$, name = R3 ] ++
       (-2,0) node [ font = \footnotesize\ttfamily ] at (R3) {10k}
       node [ left ] {$v_{i_2}$};
      \draw ([xshift = -.5cm]A.-) node [ above right ] {$v_n$}
       to [ short, *- ] ++ (0,1)
       to [ european resistor, l = $R_2$, name = R2 ] ++ (3,0) coordinate (R2r)
       node [ font = \footnotesize\ttfamily ] at (R2) {30k}
       to [ short, -* ] (R2r |- A.out);
      \draw ([xshift = -.5cm]A.+) node [ below right ] {$v_p$}
       to [ short, *- ] ++ (0,-.25)
       to [ european resistor, l_ = $R_4$, name = R4 ] ++ (0,-1.25)
       node [ rotate = 90, font = \footnotesize\ttfamily ] at (R4) {10k}
       node [ rground ] {};
    \end{circuitikz}
  \end{center}
\end{problem}
\begin{solution}

\end{solution}

\begin{problem}[15 pt]% 3
  电路图如图所示，硅 BJT 三极管的 $\beta = 100$，$C_1$ 和 $C_2$ 为隔直耦合电容.
  \begin{enumerate}
    \item 求 $Q$ 点.
    \item 求三极管放大电路的电压增益 $A_{vz}$，输入电阻 $R_i$ 和输出电阻 $R_o$.
  \end{enumerate}
  \begin{center}
    \begin{circuitikz}
      \draw (0,0) node [ npn, anchor = B ] (npn) {$T$}
       to [ european resistor, l_ = $R_b$, name = Rb ] ++ (-2,0)
       node [ font = \footnotesize\ttfamily ] at (Rb) {500}
       to [ european voltage source, l_ = $V_s$ ] ++ (0,-2)
       coordinate (ground) node [ rground ] {};
      \draw (npn.C) to [ european resistor, l = $R_c$, name = Rc ] ++ (0,1.5)
       node [ rotate = 90, font = \footnotesize\ttfamily ] at (Rc) {10k}
       node [ vcc ] {$+\qty{15}\V$} (npn.E)
       to [ european resistor, l_ = $R_e$, name = Re ] ++ (0,-1.5)
       node [ rotate = 90, font = \footnotesize\ttfamily ] at (Re) {15k}
       node [ vee ] {$-\qty{15}\V$};
      \draw (npn.E) to [ short, *-, C = $C_2$ ] ++ (1.75,0) coordinate (C2)
       (C2 |- ground) node [ rground ] {} -- (C2);
      \draw (npn.C) to [ short, *-, C = $C_1$ ] ++ (2.5,0)
       node [ below left ] {$+$}
       to [ european resistor, l = $R_L$, name = RL ] ++ (0,-1.75)
       node [ rotate = 90, font = \footnotesize\ttfamily ] at (RL) {5k}
       node [ above left ] {$V_0$}(RL |- ground)
       node [ rground ] {} node [ left ] {$-$} -- (RL);
    \end{circuitikz}
  \end{center}
\end{problem}
\begin{solution}
  
\end{solution}

\newpage

\begin{problem}[12 pt]% 2
  如图电路，$D$ 为硅二极管，$D_z$ 为，且稳定电压 $V_z = \qty3\V$，
  动态电阻 $r_z = \qty{20}\ohm$, $V_\text{DD} = \qty5\V$, $R = \qty1\kohm$，
  正弦信号 $v_s = 100\sin(2\pi \times 50t)\unit{\mV}$. 求输出电压 $v_o$.
  \begin{center}
    \begin{circuitikz}[american]
      \draw (0,0) node [ below ] {$+$} to [ short, o- ] ++ (-1.5,0)
       to [ stroke diode, invert, l_ = D ] ++ (-3,0)
       to [ sV_ = $v_s$ ] ++ (0,-2)
       to [ battery1, l_ = $V_\text{DD}$] ++ (0,-1) to [ short, -o ] ++ (4.5,0)
       node [ above ] {$-$};
      \draw (-1.5,-3) to [ short, *-, stroke diode, l_ = $D_2$ ] ++ (0,1.5)
       to [ short, -*, european resistor, l_ = $R$ ] ++ (0,1.5);
      \node at (0,-1.5) {$v_o$};
    \end{circuitikz}
  \end{center}
\end{problem}
\begin{solution}
  
\end{solution}

\begin{problem}[8 pt]% 8
  电路如图所示，试用相位平衡条件判断能不能振荡. 如果能，请说明理由；
  如果不能，也请说明理由，并最少程度上修改电路使其能产生振荡.
  \begin{center}
    \begin{circuitikz}
      \draw (0,0) node [ npn, anchor = B ] (npn2) {$T_2$} (npn2.B)
       to [ short, -* ] ++ (-.5,0) to [ european resistor, l = $R_d$ ] ++
       (0,1.5) coordinate (Rdr) to [ short, -o ] ++ (3,0)
       node [ right ] {$V_\text{CC}$};
      \draw (npn2.C) to [ short, -* ] (npn2.C |- Rdr)
       node [ njfet, anchor = D ] (njfet1) at ([xshift = -.5cm]npn2.B) {$T_1$}
       (njfet1.S) to [ european resistor, l_ = $R_1$ ] ++ (0,-1) --++
       (0,-.5) coordinate (R1r) to [ short, -* ] (R1r -| npn2.E)
       node (ground) [ rground ] {} --++ (0,.5)
       to [ european resistor, l = $R_c$ ] ++ (0,1) to [ short, -* ] (npn2.E);
      \draw (npn2.E) to [ european resistor, l = $R$ ] ++ (1.5,0)
       to [ C = $C$ ] ++ (1.5,0) coordinate (capacitor1r)
       to [ short, *-*, european resistor, l_ = $R$ ] (capacitor1r |- ground) --
       (ground);
      \draw (capacitor1r) --++ (1.25,0) coordinate (capacitor2l)
       to [ short, *-, C = $C$ ] (capacitor2l |- ground) -- (ground);
      \draw (capacitor2l) --++ (1,0) --++ (0,-3) --++ (-8,0) |- (njfet1.G);
    \end{circuitikz}
  \end{center}
\end{problem}
\begin{solution}
  
\end{solution}

\newpage

\begin{problem}[15 pt]% 4
  电路如图所示，设 MOSFET 的参数为 $V_\text{TN} = \qty1\V$，
  $K_n = \qty{0.5}{\mA/\V^2}$, $\lambda = 0$.
  \begin{enumerate}
    \item 试判断场效应管类型及其工作区.
    \item 求静态工作点.
    \item 画出电路的微变等效电路，求电路的电压增益 $A_v$，
    输入电阻 $R_i$ 和输出电阻 $R_o$.
  \end{enumerate}
  \begin{center}
    \begin{circuitikz}
      \draw (0,0) node [ below ] {$+$}
       to [ short, o-, C = $C_1$ ] ++ (1.75,0) --++ (.75,0)
       node (fet) [ anchor = G, nigfete, solderdot ] {$T$} (fet.S)
       to [ european resistor, l = $R_s$, name = Rs ] ++ (0,-1)
       node [ rotate = 90, font = \footnotesize\ttfamily ] at (Rs) {1k} --++
       (0,-1) node [ rground ] (ground) {};
      \draw (fet.D) to [ european resistor, l_ = $R_d$, name = Rd ] ++ (0,2)
       node [ rotate = 90, font = \footnotesize\ttfamily ] at (Rd) {20k}
       node [ vdd ] (vdd) {$+\qty{10}\V$};
      \draw (0,-2) node [ above ] {$-$} coordinate (negative)
       to [ short, o-* ] (negative -| ground);
      \draw (1.75,-2)
       to [ short, *-*, european resistor, l_ = $R_{g_2}$, name = Rg2 ] (1.75,0)
       node [ rotate = 90, font = \footnotesize\ttfamily ] at (Rg2) {25k}
       to [ european resistor, l_ = $R_{g_1}$, name = Rg1]
       ++ (0,2.5) --++ (0,.5) coordinate (Rg1r)
       node [ rotate = 90, font = \footnotesize\ttfamily ] at (Rg1) {100k}
       to [ short, -* ] (Rg1r -| vdd);
      \draw (fet.D) to [ short, *-o, C = $C_2$ ] ++ (2.5,0)
       node [ below ] {$v_o$};
    \end{circuitikz}
  \end{center}
\end{problem}
\begin{solution}
  
\end{solution}

\begin{problem}[12 pt]% 7
  双电源互补对称电路如图所示，设最大输出功率 $P_\text{om} = \qty9\W$,
  $R_L = \qty8\ohm$, $v_i$ 为正弦波，忽略 BJT 的饱和压降. 求
  \begin{enumerate}
    \item 正负电源 $V_\text{CC}$ 的最小值.
    \item 直流电源供给的功率 $P_v = \frac2\pi \cdot \frac{V_\text{CC}^2}{R_L}$，
    根据所求 $V_\text{CC}$ 的最小值，求放大器的效率 $\eta$.
  \end{enumerate}
  \begin{center}
    \begin{circuitikz}
      \draw (0,0) node [ vee ] {$-V_\text{CC}$} --++ (0,.25)
       node [ pnp, anchor = C ] (pnp) {$T_2$} (pnp.E)
       node [ npn, anchor = E ] (npn) {$T_1$} (npn.C) --++ (0,.25)
       node [ vcc ] {$+V_\text{CC}$};
      \draw (pnp.E) to [ short, *-o ] ++ (.9,0) --++ (.6,0)
       node [ right ] {$+$} to [ european resistor, l_ = $R_L$ ] ++ (0,-1.5)
       node [ right ] {$-$} node [ rground ] {};
      \draw ([xshift = -1.5cm]pnp.E) coordinate (vi) to [ short, *-o ]
            ([xshift = -2.5cm]pnp.E) node [ left ] {$v_i$};
      \draw (npn.B) -| (vi) |- (pnp.B);
    \end{circuitikz}
  \end{center}
\end{problem}
\begin{solution}
  
\end{solution}

\newpage

\begin{problem}[14 pt]% 5
  已知图示电路，设硅三极管 BJT 的 $\beta_1 = \beta_2 = 100$.
  \begin{enumerate}
    \item 求 $T_1$, $T_2$ 的 $Q$ 点.
    \item 计算三极管小信号模型输入电阻 $r_{be_1}$, $r_{be_2}$
    和双端输入、双端输出时的差模电压增益.
  \end{enumerate}
  \begin{center}
    \begin{circuitikz}
      \draw (0,0) node [ vee ] {$V_\text{EE}-\qty{10}\V$}
       to [ european resistor, l_ = $R_e$, name = Re ] ++ (0,1)
       node [ rotate = 90, font = \footnotesize\ttfamily ] at (Re) {9k3}
       to [ short, -* ] ++ (0,.25) coordinate (base);
      \draw (base) -| (-1,2) node (npn1) [ npn, anchor = E ] {$T_1$}
       (npn1.B) to [ short, -o ] ++ (-.25,0) node [ left ] {$V_{i_1}$}
       (npn1.C) to [ european resistor, l = $R_{e_1}$, name = Re1 ] ++ (0,2)
       node [ rotate = 90, font = \footnotesize\ttfamily ] at (Re1) {10k} --++
       (1,0) coordinate (top);
      \draw (base) -| (1,2)
       node (npn2) [ npn, xscale = -1, anchor = E, reversed ]
       {\ctikzflipx{$T_2$}}
       (npn2.B) to [ short, -o ] ++ (.25,0) node [ right ] {$V_{i_2}$}
       (npn2.C) to [ european resistor, l_ = $R_{e_2}$, name = Re2 ] ++ (0,2)
       node [ rotate = 90, font = \footnotesize\ttfamily ] at (Re2) {10k}
       to [ short, -* ] ++ (-1,0);
      \draw (top) --++ (0,.25) node [ vcc ] {$V_\text{CC}+\qty{10}\V$};
      \draw (npn1.C) to [ short, *-o ] ++ (.25,0) node [ right ] {$V_{e_1}$}
            (npn2.C) to [ short, *-o ] ++ (-.25,0) node [ left ] {$V_{e_2}$};
    \end{circuitikz}
  \end{center}
\end{problem}
\begin{solution}
  
\end{solution}

\begin{problem}[12 pt]% 6
  如图所示放大电路.
  \begin{enumerate}
    \item 请判断反馈放大电路的反馈组态.
    \item 在深度负反馈条件下，求反馈系数和闭环增益表达式.
  \end{enumerate}
  \begin{center}
    \begin{circuitikz}[scale = .75]
      \ctikzset { resistors/scale = .54, sources/scale = .6 }
      \draw (0,-1.75) node [ vee ] (vss) {$-V_\text{SS}$}
        to [ short, *-*, european current source, l = $I_O$ ] (0,0);
      \draw (0,0) -- (-1,0)
       node (fet1) [ anchor = S, nigfete, solderdot ] {$T_1$} (fet1.G) --++
       (-.5,0) to [ european resistor, l_ = $R_{g_1}$] ++ (0,-1.25)
       to [ european voltage source, l_ = $v_g$ ] ++ (0,-1) node [ rground ] {};
      \draw (0,0) -- (1,0)
       node (fet2) [ anchor = S, nigfete, solderdot, xscale = -1 ]
        {\ctikzflipx{$T_2$}} (fet2.G) --++ (.25,0)
       to [ short, *-, european resistor, l_ = $R_{g_2}$ ] ++ (0,-1.5)
       node [ rground ] {};
      \draw (fet1.D) to [ short, *- ] ++ (5,0)
       node (npn) [ npn, anchor = B ] {$T_3$}
        (npn.E) to [ european resistor, l = $R_g$ ] ++ (0,-2.25) |- (vss)
        (npn.C) to [ short, *-*, european resistor, l = $R_c$ ] ++ (0,1.5)
        coordinate (Rcr);
      \draw ([xshift = .25cm]fet2.G) to [ european resistor, l_ = $R_f$ ] ++
        (1.5,0) |- (npn.C) to [ short, *-o ] ++ (1,0) coordinate (vo+);
      \draw (fet2.D) to [ short, -*, european resistor, l = $R_{d_2}$ ]
        (fet2.D |- Rcr) (fet1.D) to [ european resistor, l = $R_{d_1}$ ]
        (fet1.D |- Rcr) -- (Rcr) to [ short, -o ] ++ (1,0)
       node [ below ] {$+V_\text{DD}$};
      \draw ([xshift = 1cm]npn.E) coordinate (vo-) to [ short, o- ] ++ (0,-.25)
       node [ rground ] {};
      \path (vo-) -- node {$v_o$} (vo+);
    \end{circuitikz}
  \end{center}
\end{problem}
\begin{solution}
  
\end{solution}

\end{document}